\section[De la consultation à l'appropriation]{De la consultation à l'appropriation : la nécessité de nouveaux outils}
Le passage d'une bibliothèque numérique de première génération, sous la forme de répertoire de PDF, à une plateforme de troisième génération, s'accompagne d'une redéfinition de la posture de l'usage en le replaçant au centre. L'usager n'est plus un \enquote{lecteur-passant} passif qui consulte un document à l'écran; il aspire à devenir un acteur actif. Il aimerait être capable d'extraire, de manipuler et de s'approprier la matière la matière textuelle et bibliographique nourrissant ces propres travaux. 

Pour guider la conception de la nouvelle interface, il convient d'adopter la théorie formulé en 2017 par Antoine Courtin. Il rappelle que: \enquote{les pratiques des chercheurs, en termes de données, se limitent bien souvent à copier des données dans un document de traitement de texte}. Face à ce constat empirique, vouloir imposers des interfaces captives ou des formats complexes et hermétiquements aux méthodes de recherche des chercheurs est inutile. La bibliothèque numérique doit proposer des services permettant l'utilisation brute des données ainsi que des services d'export s'intégrant naturellement dans les chaînes de travail personnelles des chercheurs.

Cette appropriation de la donnée par les chercheurs se heurte aujourd'hui à la volativité des sessions de recherche. Dans CujasNum, l'absence de compte utilisateur permanent condamnait le chercheur à voir son panier de documents sélectionnés disparaître dès la fermeture du navigateur. Ce mode de fonctionnement oblige le chercheur à refaire inlassablement les mêmes requêtes d'une session à l'autre afin de retrouver les documents consultés pour leurs travaux. Pour surmonter cet obstacle, la théorisation de l'espace personnel par Elina Leblanc s'avère la solution idéale pour la nouvelle bibliothèque numérique, répondant ainsi aux attentes des usagers. Leblanc préconise la  création d'un espace personnel conçu comme un véritable \enquote{carnet de lecture}. L'accès permanent à ce compte personnel ayant un historique des consultations et des recherches sauvegardées agit alors : \enquote{comme une barrière à l'oubli, mais aussi comme une économie de temps, car cela évite de recommencer ses recherches}. 

Les enquêtes qualitatives réalisées auprès des historiens du droit de Cujas confirment cette analyse. Un des professeurs explique que lorsqu'il identifie une ressource intéressante, son premier réflexe est de la télécharger afin de constituer \enquote{sa propre bibliothèque} organisée en arborescence locale sur son ordinateur. Il exprime le souhait de pouvoir transposer cette habitude de recherche à l'intérieur de la bibliothèque numérique grâce à un espace personnel persistant. De même, qu'une chercheuse insiste sur le confort d'un tel espace permettant de \enquote{mettre en favoris, lui évitant d'égarer des sources patrimoniales sur lesquelles elle doit revenir plusieurs fois au cours d'un même travail}.

Enfin, cette démarche d'autonomisation exige une pluralité de formats de sortie. Comme le souligne Elina Leblanc, proposer un éventail diversifié de formats (PDG, JPEG, exports structurés) permet de couvrir des besoins d'analyse scientifique variés dépassant le simple cadre de l'interface publique de la bibliothèque numérique de l'institution. Si le PDF conserve la préférence des chercheurs interrogés pour son caractère stable, facilement manipulable et connu, l'exportation des données bibliographiques vers des logiciels de gestion de références, comme Zotero, est devenue une attente incontournable. En plus, des formats de sortie, les pratiques de recherche des historiens révèlent le confort que représenterait la possibilité d'avoir le choix entre plusieurs options de téléchargements. Les chercheurs demandent de pouvoir faire un téléchargement partiel et ciblé. Cette demande souligne le besoin de ne part avoir à télécharger l'ensemble, lorsqu'uniquement une section ou une partie de l'ouvrage les intéresse, tout en conservant si possible la page de titre du document afin de savoir d'où provient l'extrait.

L'intégration d'un bouton d'export vers Zotero ou de la référence bibliographique intégrale s'impose comme le compagnon quotidien de l'écriture scientifique. L'affichage en clair d'une référence bibliographique normalisée \enquote{prête à être copier-coller} garantit un gain de temps précieux lors de la rédaction. Pour les utilisateurs agguerris, l'implémentation d'un module d'exportation de métadonnées riches (formats Bib Tex, CSV ou JSON), inspiré du développement mis en place dans la bibliothèque numérique \textit{Lucienne} de l'ENS permet d'automatiser l'importation massifs de corpus. En facilitant, le passage de la consultation à la possibilité de manipuler les données brutes, offre à la BIU Cujas une meilleure exploitation de ses données par les scientifiques. Il est donc nécessaire de mettre en place des accompagnements leur permettant de prendre en main ces outils.